\section{Evaluation of Results} \label{sec:evaluation}

This section presents the findings from the two experiments detailed in Section \bfref{sec:experiment-design}. The quantitative and qualitative analyses are discussed, addressing the research questions regarding content fidelity, error handling, question format generation, and alignment with Bloom's Taxonomy.

\subsection{First Experiment}
\label{sec:first-experiment}

The first experiment of the thesis focused on evaluating the content fidelity of questions generated by different \acp{llm} using two distinct prompt strategies: a common prompt and a complex prompt. In this experiment, the goal was to determine whether the complexity of the prompt influences the content fidelity and final scoring of the generated questions, addressed in the first research question:

\requ{1}{To what extent do \acp{llm} adhere to the content of diverse provided instructional texts when generating questions?} Experiments 1a and 1b follow both quantitative and qualitative approaches each to evaluate the generated questions properly. The quantitative analysis solely focuses on the analysis of content fidelity, while the qualitative analysis assesses the quality of the generated questions based on a certain rubric. This rubric also covers the question's \textit{Correctness}, which provides an underlying criterion for assessing content fidelity on expert-level.

\paragraph*{Quantitative Analysis} 

Two metrics were employed to assess content adherence: the calculation of \textit{Semantic Similarity} using cosine similarity between resulting semantic vectors -- one for the source material file, and one for the question file -- from the embedding model, and \ac{llm}-based \textit{Adherence Score} analysis between these files as an alternate approach to evaluate content fidelity.

\object{Experiment 1a -- Content Fidelity} The manner of evaluation revealed several key patterns in the \acp{llm}' behavior. The question type for each question was not explicitly constrained, but to provide a valid analysis -- especially for the experts --, the questions' output formats were limited to either Open-Ended or Multiple-Choice questions. This resulted in a trend toward \acp{mcq}, with only 9 out of 168 questions being Open-Ended in Experiment 1a -- with 7 questions from OpenAI, and 2 from Google --, based on the complex prompt. With regard to this finding, the trend was consistent across the \ac{llm} set, demonstrating a bias toward structured question formats.

As an introductory step, to observe if both approaches are valid, the \acp{llm} were prompted to generate a question for each ISO-OSI layer without using an input source. Here, a key observation is that the cosine similarity approach did not reliably distinguish between questions generated with or without input sources. As shown in Figure \bfref{fig:exp1a-source-scores}, the scores for the ones without a source -- maintaining a mean $\approx 0.533$ and a median $\approx 0.538$ -- were quite similar to those using actual source material with means between $[0.496_{\text{transcript}},0.637_{\text{script}}]$ and nearby medians at $[0.511_{\text{transcript}},0.651_{\text{script}}]$, indicating that cosine similarity is not a robust indicator of content fidelity for these kinds of texts. 

\begin{figure}[htbp]
    \centering
    \includegraphics[width=.9\textwidth]{../../../40_evaluation/exp1/quantitative/plots/exp1a_no_source/input_source_vs_no_source_comparison.png}
    \caption{Quantitative Distributions for Input Sources vs.\ Run without any Source (Exp1a).}
    \label{fig:exp1a-source-scores}
\end{figure}

In contrast, the \textit{Adherence Score} provided a clearer differentiation. Using no source led to noticeably lower adherence overall in the mean $\approx 0.714$, the median $\approx 0.768$ and a broader range of scores. Additionally, a certain trend in cosine similarity across the sources is visible in Figures \bfref{fig:exp1a-source-scores} and \bfref{fig:exp1a-llm-source-comparison}. There is to be seen that the precise lecture script is indicating best results, Tanenbaum excerpts are situated in the middle, and the extensive transcript parts equipped with examples keep the widest range of scores, the lowest mean and median values. This trend appears to be influenced by the amount of words in the source materials rather than content quality itself, limiting its overall usefulness. Moreover, the \textit{Adherence Score} showed few outliers in the box plot, and the scoring in the heatmap shows quite individual \ac{llm} performance according to the diverse input sources, overall located above 90\%.

\begin{figure}[htbp]
    \centering
    \includegraphics[width=.8\textwidth]{../../../40_evaluation/exp1/quantitative/plots/exp1a/combined_llm_vs_input_source.png}
    \caption{Quantitative Heatmaps for LLM vs.\ Input Source Performance (Exp1a).}
    \label{fig:exp1a-llm-source-comparison}
\end{figure}

Regarding the \ac{llm}-sorted observations, Figure \bfref{fig:exp1a-llm-scores} reveals that the cosine similarity scores were lower and also differ a lot from the \textit{Adherence Scores}. Expert results by the primary supervisor, later discussed under \textit{Correctness} in e.g.\ Figure \bfref{fig:exp1a-correctness-heatmaps}, show that cosine similarity was unreliable for this task, as this did not properly reflect the content fidelity of the questions in Experiment 1a. The \textit{Adherence Scores} in this quantitative analysis, however, provided a quite nuanced view of model performance, showing that DeepSeek and OpenAI were more consistent in adhering to the source content, while Anthropic and Google exhibited slightly lower and more variable scores.

\begin{figure}[htbp]
    \centering
    \includegraphics[width=.9\textwidth]{../../../40_evaluation/exp1/quantitative/plots/exp1a/dist_by_llm.png}
    \caption{Quantitative Scoring Distributions by LLM (Exp1a).}
    \label{fig:exp1a-llm-scores}
\end{figure}

To get a proper view of the content adherence, the evaluation also included a comparison of the performance of the \acp{llm} against the used prompt type, as shown in Figure \bfref{fig:exp1a-llm-prompt-comparison}. Here, the results between cosine similarity and \textit{Adherence Scores} were again significantly different. On the one hand, the hardly comprehensible cosine similarity tends to higher scoring values for the complex prompt. On the other hand, \textit{Adherence Scores} show that the common prompt leads to better results -- best for OpenAI ($\approx 0.992$ vs.\ $0.94$) and DeepSeek ($\approx 0.973$ vs.\ $0.945$) --, while the complex prompt leads to slightly lower scores instead for all models, especially for Google and Anthropic. This is also indicated by the expert results by the primary supervisor, which indicate that the common prompt is more effective in content adherence than the complex prompt.

\begin{figure}[htbp]
    \centering
    \includegraphics[width=.8\textwidth]{../../../40_evaluation/exp1/quantitative/plots/exp1a/combined_llm_vs_prompt_type.png}
    \caption{Quantitative Heatmaps for LLM vs.\ Prompt Type Performance (Exp1a).}
    \label{fig:exp1a-llm-prompt-comparison}
\end{figure}

These findings suggest that, while cosine similarity may be suitable for short, sentence-level comparisons -- as the \textit{Sentence Transformer} expression suggests --, it is less appropriate for evaluating content fidelity between generated questions and also long, more complex educational texts. Consistent context e.g.\ up to 30 words might improve cosine similarity's reliability, but in this setting, it will be disregarded as a primary metric for Experiment 1a, since the thesis deals with texts which are longer than a single sentence and contain complex information. As an example, \cite{li_planning_2024} used BERTScore for semantic similarity scoring, but between generated and ground truth question-answer pairs. This is not the case in this thesis, since the source texts and \ac{llm} outputs differ significantly.

\pagebreak

\object{Experiment 1b -- Error Handling} This sub-experiment focused on evaluating how well the \acp{llm} handle manipulated lecture script content to test their ability to detect errors and what the models generated based on that content. A total of 4/56 resulted in open-ended questions, one time DeepSeek using the common prompt, two times Google plus one time OpenAI both using the complex prompt. Overall, it is worth noting that using the cosine similarity approach here -- as to be seen in Figure \bfref{fig:exp1b-llm-scores} -- is less useful than in Experiment 1a, since the embedding model can not detect if a question by an \ac{llm} follows the manipulated content flawlessly, or if the model detects the manipulation and e.g.\ generated the question inversely.

\begin{figure}[htbp]
    \centering
    \includegraphics[width=.9\textwidth]{../../../40_evaluation/exp1/quantitative/plots/exp1b/dist_by_llm.png}
    \caption{Quantitative Scoring Distributions by LLM (Exp1b).}
    \label{fig:exp1b-llm-scores}
\end{figure}

\pagebreak

\noindent Some examples -- translated into English -- of such questions are:

\begin{itemize}
    \item \enquote{Which of the following statements describes a \textbf{real core task} of the Data Link Layer in the ISO-OSI model, based on \textbf{correct} technical knowledge?} or 
    \item \enquote{Which property is \textbf{NOT} assigned to the Transport Layer according to the \textbf{description?}}.
\end{itemize}

However, the \textit{Adherence Scores} illustrate a proper course in the plots, showing that Google and Anthropic maintain more often to the manipulated content -- based on value range and the high means $\approx 0.92$ and $\approx 0.904$. Instead, OpenAI and especially DeepSeek detected the manipulation in some cases and generated questions where this aspect is recognizable. The rightly present outliers for OpenAI separate its mean from its high median by a value of $0.988-0.806=0.182$. Moreover, the DeepSeek model -- the oldest of the four chosen models -- convinced regarding the lowest median $\approx 0.88$, mean adherence $\approx 0.731$ and a wide range of values, indicating that it often recognized the OSI-layer manipulation.

Furthermore, the used type of prompt had an impact on the results, as shown in Figure \bfref{fig:exp1b-llm-prompt-comparison}. In this case, both cosine similarity and \textit{Adherence Scores} indicate that the common prompt led to higher results than the complex prompt, while lower results are the goal in this sub-experiment instead. As stated, the cosine similarity cannot reliably determine if an \ac{llm} follows the manipulated content or if it detects it and generates a question based on that text. In contrast to that, the \textit{Adherence Scores} provide more reliable results, showing that the complex prompt led to lower scores for all \acp{llm}. Like in Experiment 1a, DeepSeek and OpenAI were more convincing in their performance -- here with lower mean \textit{Adherence Scores} ($\approx 0.598$ and $\approx 0.619$) --, while Anthropic and Google still have shown higher scores. Overall, OpenAI and DeepSeek were the models that recognized the manipulation more often, but in this sub-experiment, the complex prompt led to \enquote{better} \textit{Adherence Scores} than the common prompt.

\begin{figure}[htbp]
    \centering
    \includegraphics[width=.8\textwidth]{../../../40_evaluation/exp1/quantitative/plots/exp1b/combined_llm_vs_prompt_type.png}
    \caption{Quantitative Heatmaps for LLM vs.\ Prompt Type Performance (Exp1b).}
    \label{fig:exp1b-llm-prompt-comparison}
\end{figure}

\pagebreak

\paragraph*{Qualitative Analysis -- Supervisor-based} 

To evaluate the generated questions properly and expand the quantitative analysis, a sampled qualitative analysis of Experiment 1a was conducted. This involved a detailed rubric evaluation by the primary thesis supervisor and four department staff members.

\object{Experiment 1a -- Content Fidelity} The following criteria were used to assess the questions: \textit{Relevance}, \textit{Clarity}, \textit{Answerability}, \textit{Challenging}, \textit{Value}, \textit{Language}, and \textit{Correctness}. Each criterion was rated on a scale from 1 to 10, with 10 being the best score. The experts were asked to evaluate the questions based on these criteria, and the results were aggregated to provide an overall assessment of the question quality. In this paragraph, the supervisor's results are firstly discussed, while the staff members' results are presented right after. 

The supervisor -- a professor with extensive lecturing experience -- analyzed the questions first, and the staff members were then asked to evaluate the same questions based on the same criteria. In the meantime, the supervisor noticed that the used rubric was not sufficient to cover all qualitative aspects of the questions, leading to an extension of the rubric. By this, the mentioned criteria \textit{Value} and \textit{Language} were added to the rubric, which were not part of the firstly used rubric. The novel \textit{Value} criterion was introduced to rate the meaningfulness of the questions in an e.g.\ educational context, while the \textit{Language} criterion focused on the linguistic complexity and precision of the questions -- since \textit{Clarity} focused the content clarity and logical comprehension.

Regarding Figure \bfref{fig:exp1a-llm-analysis}, certain patterns emerged from the supervisor-based evaluation. Most questions were rated highly \textit{relevant} to the material's topic, with OpenAI and Google as the highest scores. While every model has a median of 10, only the OpenAI model could achieve a mean of 10 as well at \textit{Relevance}, also to be seen at the \textit{Clarity} criterion next to this criterion. Here, DeepSeek was also rated high as a valuable competitor of OpenAI. In contrast, Anthropic and Google both had a wide scale of points and means of 7.33 and 7.64, indicating that these two models were more inconsistent at generating clear and logically understandable questions. 

Moreover, the \textit{Answerability} criterion -- based on the source text -- has shown that OpenAI and DeepSeek were the best-performing models again (with means of 9.33 and 8.5), while Anthropic's and Google's scores (with means of 7.5 and 7.45) were spreading more to lower scores instead. Contrary to this, the questions of the latter two models were rated as more challenging -- with Anthropic being most consistent with a mean of 6.67 --, compared to the significantly lower OpenAI and DeepSeek scores in both mean (3.17 and 3.5) and median (both at 2) in this regard.

The educational \textit{Value} criterion indicates that the OpenAI questions were the most valuable ones with a mean and median of 10. The other models were inconsistent in this regard, with noticeably wide ranges of scores, means between 6 and 7 and DeepSeek having the lowest median of 7.

The \textit{Language} criterion has a similar trend like \textit{Clarity}, with OpenAI and DeepSeek performing best with a mean of 10 each. Anthropic and Google were rated lower and more spread across the scale -- resulting in less means of 8.25 and 7.67, showing that these models generated questions which were less precise and unnecessarily complex. 

\pagebreak

\begin{figure}[H]
    \centering
    \includegraphics[width=.9\textwidth]{../../../40_evaluation/exp1/qualitative/plots/supervisor_exp1a_llm_analysis.png}
    \caption{Supervisor-based LLM Scores across every proposed Criterion (Exp1a).}
    \label{fig:exp1a-llm-analysis}
\end{figure}

\pagebreak

Ultimately, the \textit{Correctness} criterion, to assess the overall content fidelity of the questions, has shown that OpenAI was the best performing model again with a 10/10 mean. DeepSeek and Anthropic are situated right after, with an exemplary question by DeepSeek in Appendix \bfref{lst:1a-supervisor-good-example}. The questions generated by Google were rated all across the board, which shows that the model struggled generating correct questions which shall be solely based on the source material.

Overall, the \textit{Total Score} shows that OpenAI achieved the highest scores with a mean of 62.5/70, a pretty similar median of 62 and a very low range of values. Followed by DeepSeek and Anthropic, the means are situated lower at 53.25 and 52.83 -- with more spread values --, while Google was rated lowest with a mean of 46.17 and most spread scores at all.

Since the \textit{Correctness} criterion is the most important one for Experiment 1a, the following Heatmaps in Figure \bfref{fig:exp1a-correctness-heatmaps} also show certain patterns to be observed. OpenAI still achieved 10/10 points everywhere in both heatmaps, which is a surprising contrast to the other models. Google's correctness has become worse with decreasing complexity of the input source, falling from 10/10 at transcript to 3.3/10 using the precise script. A somewhat similar course can be seen for DeepSeek, but less pronounced. There, the model fell from 10/10 at transcript to 7.5/10 at the script, with also more stable standard deviations. Anthropic's correctness scores were more stable, scaling from 8.8/10 at (tran-)script to 7.5/10 at the Tanenbaum book excerpts.

\begin{figure}[H]
    \centering
    \includegraphics[width=.8\textwidth]{../../../40_evaluation/exp1/qualitative/plots/supervisor_exp1a_correctness_heatmaps.png}
    \caption{Supervisor-based Correctness Heatmaps for different LLM Combinations (Exp1a).}
    \label{fig:exp1a-correctness-heatmaps}
\end{figure}

Regarding the heatmap focusing on the prompt type, there is a trend that the common prompt leads to higher correctness scores than the complex prompt, which was also depicted via the quantitative \textit{Adherence Score} in Figure \bfref{fig:exp1a-llm-prompt-comparison}. OpenAI and DeepSeek achieved 10/10 points at the common prompt, Anthropic is at 9.2 points and Google at 8.3 points. The complex prompt instead had decreased correctness scores, except for OpenAI. DeepSeek and Anthropic were both at 7.5/10 points, while Google solely achieved 3.3 points with high standard deviations for both prompt types.

\pagebreak

Since the input source and the prompt type are also worth discussing, the results of the qualitative analysis are summarized in Tables \bfref{tab:exp1a-source-results} and \bfref{tab:exp1a-prompt-results}. The tables show the mean, standard deviation, and median for each criterion. First of all, visible in Table \bfref{tab:exp1a-source-results}, almost every criterion has a median of 10/10 points, except for \textit{Challenging}. The mean values differ in contrast to the median. 

Overall, there is a trend that the extensive transcript excerpts lead to the best results -- especially in means of \textit{Clarity} (9.33), \textit{Answerability} (9.33), and \textit{Correctness} (9.67) --, followed by the precise script and the Tanenbaum book excerpts. The latter two sources were on par at the \textit{Language} criterion with a mean value of 8.88. In contrast, the questions based on the script were most challenging with a mean of 5.5, followed by Tanenbaum here with 5.38. The transcript-based questions were rated lowest in \textit{Challenging}, only achieving a mean of 2.93 points. Furthermore, while the transcript questions were the most accurate regarding the given context to the \acp{llm} with a mean of 9.67, the Tanenbaum's (8.33) and especially script's \textit{Correctness} (7.67) scores were rated lower instead.

Even if the transcript appears to be the best source material for many criteria, the script's mean \textit{Total Score} was the highest with 55.06/70 points, followed closely by the transcript with 54.75 points and the Tanenbaum excerpts with 51.25 points. Regarding the median values, the transcript was rated highest with 62/70 points instead -- in contrast to the script (61) and Tanenbaum (61.5) --, but the questions generated based on the transcript lose points due to not being as challenging as the other sources.

\begin{table}[htbp]
    \centering
    \caption[Supervisor-based qualitative Analysis Results by Input Source (Exp1a).]{Supervisor-based qualitative Analysis Results by Input Source (Exp1a). The best means are highlighted in \textbf{bold}, while the best medians are \textit{italicized}.}
    \label{tab:exp1a-source-results}
    \resizebox{\textwidth}{!}{%
    \begin{tabular}{lccccccccc}
        \toprule
        \textbf{Source} & \textbf{Metric} & \textbf{Relevance} & \textbf{Clarity} & \textbf{Answerability} & \textbf{Challenging} & \textbf{Value} & \textbf{Language} & \textbf{Correctness} & \textbf{Total Score} \\
        \midrule
        \multirow{3}{*}{Script} 
        & Mean & 9.38 & 8.50 & 8.25 & \textbf{5.50} & 7.38 & 8.88 & 7.67 & \textbf{55.06} \\
        & Std & 2.03 & 2.37 & 2.52 & 3.54 & 3.32 & 2.19 & 3.72 & 10.23 \\
        & Median & 10.0 & 10.0 & 10.0 & \textit{6.0} & 10.0 & 10.0 & 10.0 & 61.0 \\
        \midrule
        \multirow{3}{*}{Tanenbaum} 
        & Mean & 9.12 & 8.25 & 7.12 & 5.38 & 7.14 & 8.88 & 8.33 & 51.25 \\
        & Std & 2.19 & 3.26 & 3.79 & 4.05 & 4.05 & 2.19 & 3.26 & 14.88 \\
        & Median & 10.0 & 10.0 & 10.0 & 5.0 & 10.0 & 10.0 & 10.0 & 61.5 \\
        \midrule
        \multirow{3}{*}{Transcript} 
        & Mean & \textbf{9.47} & \textbf{9.33} & \textbf{9.33} & 2.93 & \textbf{7.87} & \textbf{9.19} & \textbf{9.67} & 54.75 \\
        & Std & 2.07 & 2.09 & 1.80 & 2.49 & 4.10 & 1.94 & 1.29 & 15.88 \\
        & Median & 10.0 & 10.0 & 10.0 & 2.0 & 10.0 & 10.0 & 10.0 & \textit{62.0} \\
        \bottomrule
    \end{tabular}}
\end{table}

Besides the input sources, the results in Table \bfref{tab:exp1a-prompt-results} show a trend, which was depicted in the quantitative analysis as well. The common prompt led to higher scores in almost every criterion, while the complex prompt was rated lower. The only criterion where the complex prompt was rated higher was \textit{Challenging} in both mean and median values (6.87 vs.\ 2.50 and 8.0 vs 2.0). Overall, the common prompt outperformed the complex prompt with a noticeable gap in the mean \textit{Total Score} of 60.12 vs.\ 47.25 points, as well as in the median of 62.0 vs.\ 48.5 points. 

The ratings by the supervisor indicate the following: the complex prompt led to less relevant questions, which were less clear and especially less answerable, did not share the same value by far as the common prompt questions, and the questions were overly complex and less precisely formulated. The \textit{Correctness} to the source material was also rated lower for the complex prompt, which is crucial for this experiment, but the questions were noticeably more challenging instead. Overall, the transcript-based questions were rated highest, but the lack of challenge in these questions led to a lower total score compared to the script-based questions. The questions created by OpenAI were rated highest in many criteria, while the questions were less challenging. This observation was also reflected in DeepSeek's results, which were also rated high, but struggling with the \textit{Challenging} and \textit{Value} criteria. Anthropic and Google were more challenging, but their questions overall lacked \textit{Clarity}, \textit{Answerability}, \textit{Language}, and \textit{Correctness}.

\begin{table}[htbp]
    \centering
    \caption[Supervisor-based qualitative Analysis Results by Prompt Type (Exp1a).]{Supervisor-based qualitative Analysis Results by Prompt Type (Exp1a).}
    \label{tab:exp1a-prompt-results}
    \resizebox{\textwidth}{!}{%
    \begin{tabular}{lccccccccc}
        \toprule
        \textbf{Prompt} & \textbf{Metric} & \textbf{Relevance} & \textbf{Clarity} & \textbf{Answerability} & \textbf{Challenging} & \textbf{Value} & \textbf{Language} & \textbf{Correctness} & \textbf{Total Score} \\
        \midrule
        \multirow{3}{*}{Common} 
        & Mean & \textbf{10.00} & \textbf{10.00} & \textbf{9.92} & 2.50 & \textbf{9.55} & \textbf{10.00} & \textbf{9.35} & \textbf{60.12} \\
        & Std & 0.00 & 0.00 & 0.41 & 1.69 & 2.13 & 0.00 & 2.29 & 6.56 \\
        & Median & 10.0 & 10.0 & \textit{10.0} & 2.0 & \textit{10.0} & 10.0 & 10.0 & \textit{62.0} \\
        \midrule
        \multirow{3}{*}{Complex} 
        & Mean & 8.61 & 7.30 & 6.43 & \textbf{6.87} & 5.48 & 7.96 & 7.63 & 47.25 \\
        & Std & 2.79 & 3.23 & 3.36 & 3.66 & 3.92 & 2.56 & 3.48 & 15.93 \\
        & Median & 10.0 & 10.0 & 6.0 & \textit{8.0} & 4.0 & 10.0 & 10.0 & 48.5 \\
        \bottomrule
    \end{tabular}}
\end{table}

\object{Experiment 1b -- Error Handling} A second run was conducted to evaluate \ac{llm}-based questions based on manipulated lecture script content. The qualitative analysis of Experiment 1b was done using the same rubric, while \textit{Correctness} was substituted by \textit{Manipulation Handling}. This was introduced to assess how well the \acp{llm} handled the manipulated lecture script content. The same five experts were asked to evaluate these questions. First, the thesis supervisor analyzed the questions, followed by the staff members regarding the \bfhyperref{sec:exp1b-staff-analysis}{next} paragraph.

Unfortunately, the supervisor's evaluation in Experiment 1b revealed that the rubric was not sufficient for this assessment, because the contextual material is manipulated. Based on this, the \textit{Manipulation Handling} criterion was solely analyzed, and the results are presented in Table \bfref{tab:exp1b-manipulation-handling-supervisor}. If this issue was noticed earlier, the questions could have been evaluated by solely this single criterion, resulting in creating a bigger sampling set, where e.g.\ the full set of questions could have been analyzed. However, the sampled set of questions is too small to draw any significant conclusions, since certain questions where the models noticed the manipulation were not rated by the experts.

\begin{table}[htbp]
    \centering
    \caption{Supervisor-based Manipulation Handling Analysis by LLM and Prompt Type (Exp1b).}
    \label{tab:exp1b-manipulation-handling-supervisor}
    \resizebox{.5\textwidth}{!}{%
    \begin{tabular}{llccc}
        \toprule
        & & \multicolumn{3}{c}{\textbf{Manipulation Handling}} \\
        \cmidrule{3-5}
        \textbf{LLM} & \textbf{Prompt Type} & \textbf{Mean} & \textbf{Std} & \textbf{Median} \\
        \midrule
        \multirow{2}{*}{Anthropic} 
        & Common & 0.0 & 0.0 & 0.0 \\
        & Complex & \textbf{2.5} & 3.54 & \textit{2.5} \\
        \midrule
        \multirow{2}{*}{DeepSeek} 
        & Common & 2.5 & 3.54 & 2.5 \\
        & Complex & 2.5 & 3.54 & 2.5 \\
        \midrule
        \multirow{2}{*}{Google} 
        & Common & 0.0 & 0.0 & 0.0 \\
        & Complex & \textbf{2.5} & 3.54 & \textit{2.5} \\
        \midrule
        \multirow{2}{*}{OpenAI} 
        & Common & 0.0 & 0.0 & 0.0 \\
        & Complex & 0.0 & --- & 0.0 \\
        \bottomrule
    \end{tabular}}
\end{table}

\pagebreak

In this table, there are results based on 16/56 questions, two for each prompt type and \ac{llm}. The limited sample size led to several observations. The \textit{Manipulation Handling} criterion was rated with a mean and median of 2.5/10 or 0, commonly indicating that the models did not handle the manipulation well. The values of 2.5 illustrate that one of the two questions was rated with 5/10 points, while the other one was rated with 0/10 points. 5/10 points show that the systems still followed the manipulated content, but in the formulation of the questions, there was a hint that the manipulation was noticed.

Anthropic and Google were rated with a mean of 2.5/10 for the complex prompt each, with Google's 5/10 question slightly distancing from the manipulated script (Appendix \bfref{lst:1b-supervisor-good-example}). The questions based on the common prompt were rated with 0/10 points for both models, with Google providing a proper example of following the manipulation (Appendix \bfref{lst:1b-supervisor-bad-example}). DeepSeek's questions were rated with 2.5/10 for both prompt types, while OpenAI's questions show 0/10 points for each prompt type. The latter is a surprising result, because in the quantitative analysis, OpenAI was performing well with the manipulation, as well as DeepSeek -- both providing few essentially low \textit{Adherence Scores}. One of the two complex prompt questions the OpenAI model generated could not be rated by the supervisor, since it was hard to evaluate its manipulation handling here. This question (Appendix \bfref{lst:1b-supervisor-nan-example}) results in a NaN value -- which is why the standard deviation for OpenAI's complex prompt questions is not available --, showing again that the sampling size was too small.

Overall, the supervisor-based evaluation does not provide significant insights into the \acp{llm}' \textit{Manipulation Handling}, but the models particularly struggled with this task. The sampled question sets were set as-is, because rating a more vast number of questions -- focusing on the question's context, and a comprehensive rubric -- was not feasible within the raters' time constraints.

\object{Supervisor-based Analysis Complications} During this evaluation, several complications occurred. Firstly, certain questions could not be fully assessed for specific criteria, because these questions were not formulated properly, leading to NaN values, such as Appendix \bfref{lst:1a-supervisor-bad-nan-example}. This issue was present in both Experiments 1a and 1b, possibly affecting the calculations for the evaluation.

A closer examination of the questions revealed issues for the models. Regarding Anthropic, several questions were problematic, especially in terms of \textit{Challenging}, \textit{Value} and \textit{Correctness}, often coupled with the complex prompt used. In some cases, answer options were incorrect or did not align with the expected content. DeepSeek frequently mentioned \enquote{according to the text} in their questions, which also was a common issue at Google's questions. These were often formulated far too broadly at Google, with some containing answer options which were factually incorrect, coupled with unclear phrasing. OpenAI, in contrast, showed occasionally the \enquote{according to the text} phrasing, but overall, the questions were fine but less challenging, e.g.\ comparable to DeepSeek's questions.

Since criteria like \textit{Answerability} could not be reliably applied to the manipulated content in Experiment 1b, the evaluation was limited to the \textit{Manipulation Handling} criterion. This was done because the rubric's criteria were comprehensively assessed in Experiment 1a, and the manipulation handling was the only criterion that could be reliably applied to the manipulated lecture script.

\vp

% --------------------------------------------------------------------------------

\paragraph*{Qualitative Analysis -- Staff-based Comparison} \label{sec:exp1b-staff-analysis}

In this sub-experiment, four staff members were asked to evaluate the same question set, as the supervisor has done before. This paragraph discusses the staff members' outcomes primarily including the similarities and differences to the supervisor's results.

\object{Experiment 1a -- Content Fidelity} The staff members' evaluation revealed several notable similarities and differences compared to the supervisor's assessment. As shown in Figure \bfref{fig:exp1a-llm-analysis-staff}, the overall trends remained consistent, but the staff members were more stringent in their scoring.

A key similarity was the identification of OpenAI and DeepSeek as the leading models in most criteria such as \textit{Clarity}, \textit{Answerability}, and \textit{Language}. However, where the supervisor saw a clear advantage for OpenAI (e.g. mean of 10 in \textit{Relevance}), the staff rated the models closer together, with OpenAI (9.12) only slightly ahead of the others. This suggests that the performance gap between the models is mainly perceived as smaller by the staff.

As depicted by the supervisor, the \textit{Challenging} criterion also remained the lowest-scoring criterion across all models, with means below 5 for all \acp{llm} at the staff-based one. DeepSeek (2.92) and OpenAI (3.35) performed lowest again, while Anthropic (4.65) and Google (4.81) were rated as slightly more challenging, but not as significantly as in the supervisor's evaluation. Especially the ratings for Anthropic and Google were different from the supervisor's, showing that the staff members found these models' questions less challenging than the supervisor did instead.

Furthermore, the staff's more critical assessment is visible in the linguistic evaluation. For the \textit{Language} criterion, they rated Anthropic (7.4) and Google (6.96) lower than the supervisor did (8.25 and 7.67 respectively). Moreover, OpenAI and DeepSeek were also rated lower than the supervisor's 10/10 scores for both models, with means of 8.6 each. These values differ from the supervisor's evaluation, but still show that the staff members found these models' questions more precise instead of linguistically complex. The \textit{Value} criterion also slightly differed from the supervisor's evaluation, with all means below 8 points. The main difference was that the staff members rated OpenAI with a mean of 6.48 instead of 10/10 points, like the supervisor did. As a negative example for these two criteria, in Appendix \bfref{lst:1a-staff-bad-example}, a question by Google is shown, formulated inappropriately by using a negation in the question itself. 

For \textit{Correctness}, the staff-based evaluation confirmed OpenAI's high performance (8.42) with DeepSeek beyond (8.48) -- both with narrow value ranges --, though without the clear separation the supervisor observed. Close behind these models, Anthropic and Google were rated lower with means of 7.94 and 8.15. The model ratings differ from the supervisor's, since the \textit{Correctness} criterion was implemented choosing ratings of 0, 5, or 10. The thesis supervisor solely used these ratings, while the staff members rated values in between. The supervisor could have rated the questions differently, not rating this many questions with 10/10 points, but rather values between 0 and 5, or 5 and 10, which would have led to a more similar evaluation to the staff members' one.

\pagebreak

\begin{figure}[H]
   \centering
    \includegraphics[width=.9\textwidth]{../../../40_evaluation/exp1/qualitative/plots/staff_exp1a_llm_analysis.png}
    \caption{Staff-based LLM Scores across every proposed Criterion (Exp1a).}
    \label{fig:exp1a-llm-analysis-staff} 
\end{figure}

\pagebreak

Overall, the \textit{Total Scores} revealed that while OpenAI (mean of 53.5) remained the top performer, but less pronounced compared to the supervisor's clear preference (62.5 vs. DeepSeek with 53.25). The results for the other models were more similar to each other, with means around 50/70 points. This is a contrast to the supervisor's evaluation, where Google especially was often rated lower. 

Regarding Figure \bfref{fig:exp1a-correctness-heatmaps-staff}, the \textit{Correctness} scores focusing on the input sources differ from the supervisor's evaluation. There is a minor trend similar to the supervisor's evaluation, where OpenAI and DeepSeek tend to be rated higher than the other models, also visible in the \textit{Adherence Score}. Contrary to the supervisor, the staff members rated OpenAI's questions based on Tanenbaum with a mean of 7.8, which is significantly lower than the supervisor's 10/10 points. Also, the trend that the transcript-based questions were rated highest is not visible in the staff members' mixed picture.

In the second staff heatmap, the prompt type had the same impact on the correctness scores, with the common prompt leading to higher scores than the complex one. The models were on par with values all above 9/10 in the mean. In the supervisor's evaluation, Google was rated with 8.3 in the common prompt instead. Focusing the complex prompt, the supervisor also rated Google significantly lower with a mean of 3.3/10, while the staff members did with 6.8 points. 

\begin{figure}[htbp]
    \centering
    \includegraphics[width=.8\textwidth]{../../../40_evaluation/exp1/qualitative/plots/staff_exp1a_correctness_heatmaps.png}
    \caption{Staff-based Correctness Heatmaps for different LLM Combinations (Exp1a).}
    \label{fig:exp1a-correctness-heatmaps-staff}
\end{figure}

In the following, the results of the staff-based qualitative analysis are summarized in Table \bfref{tab:exp1a-source-results-staff} to observe how the staff members have rated based on the input source. These revealed a minor shift in source preference compared to the supervisor's findings in Table \bfref{tab:exp1a-source-results}. While the supervisor consistently favored transcript-based questions as the best source -- with the highest means in \textit{Clarity} (9.33), \textit{Answerability} (9.33), \textit{Value} (7.87), \textit{Language} (9.19), and \textit{Correctness} (9.67) -- the staff members slightly identified script-based questions as the overall best performer. The script achieved the highest staff ratings in \textit{Relevance} (9.16), \textit{Challenging} (4.06), \textit{Value} (6.67), \textit{Correctness} (8.33), and \textit{Total Score} (52.14), contrasting with the supervisor's input source preference.

\pagebreak

\begin{table}[htbp]
    \centering
    \caption[Staff-based qualitative Analysis Results by Input Source (Exp1a).]{Staff-based qualitative Analysis Results by Input Source (Exp1a). The best means are highlighted in \textbf{bold}, while the best medians are \textit{italicized}.}
    \label{tab:exp1a-source-results-staff}
    \resizebox{\textwidth}{!}{%
    \begin{tabular}{lccccccccc}
        \toprule
        \textbf{Source} & \textbf{Metric} & \textbf{Relevance} & \textbf{Clarity} & \textbf{Answerability} & \textbf{Challenging} & \textbf{Value} & \textbf{Language} & \textbf{Correctness} & \textbf{Total Score} \\
        \midrule
        \multirow{3}{*}{Script} 
        & Mean & \textbf{9.16} & 7.86 & 8.17 & \textbf{4.06} & \textbf{6.67} & 7.89 & \textbf{8.33} & \textbf{52.14} \\
        & Std & 1.64 & 2.55 & 2.65 & 2.34 & 3.18 & 2.48 & 2.56 & 10.82 \\
        & Median & 10.0 & 9.0 & 10.0 & \textit{4.0} & \textit{8.0} & 8.0 & 9.5 & \textit{55.0} \\
        \midrule
        \multirow{3}{*}{Tanenbaum} 
        & Mean & 8.77 & 7.62 & 7.66 & 3.94 & 5.77 & 7.77 & 8.17 & 49.69 \\
        & Std & 2.09 & 2.81 & 3.10 & 2.44 & 3.41 & 2.63 & 2.79 & 11.83 \\
        & Median & 10.0 & 8.0 & 9.5 & 3.0 & 6.0 & 8.0 & 10.0 & 52.0 \\
        \midrule
        \multirow{3}{*}{Transcript} 
        & Mean & 8.89 & \textbf{7.89} & \textbf{8.19} & 3.80 & 6.06 & \textbf{8.02} & 8.23 & 51.08 \\
        & Std & 2.02 & 2.85 & 2.81 & 2.32 & 3.10 & 2.72 & 2.73 & 11.23 \\
        & Median & 10.0 & \textit{10.0} & 10.0 & 3.5 & 6.0 & \textit{10.0} & 10.0 & 54.0 \\
        \bottomrule
    \end{tabular}}
\end{table}

Furthermore, the rating gaps between the two evaluation groups are notable. For instance, the supervisor rated the transcript-based questions significantly higher in \textit{Clarity} (9.33 vs. 7.89), \textit{Answerability} (9.33 vs. 8.19), and \textit{Correctness} (9.67 vs. 8.23). Similarly, the script-based questions have shown several differences in the \textit{Challenging} criterion, where the supervisor rated them at 5.5 compared to the staff's 4.06.

Despite these scoring differences, both evaluations commonly identified Tanenbaum excerpts as the weakest source material (excluding \textit{Challenging}), with the staff rating it lowest in \textit{Total Score} (49.69), and the supervisor placing it similarly (51.25). It struggled across both evaluation groups, particularly in clear and educationally valuable questions. Both exhibited the same trend in the \textit{Total Score}, where the script was rated slightly highest with a mean of 52.14 -- compared to the supervisor's 55.06 --. This is followed by the transcript with 51.08 -- compared to the supervisor's 54.75. In addition, an example of a properly formulated question by OpenAI -- based on the transcript -- is shown in Appendix \bfref{lst:1a-staff-supervisor-good-example}, rated almost perfect across the criteria, indicating that both groups found this question clear, precise, and answerable, focusing on the content of the transcript.

To further investigate the prompt differences, Table \bfref{tab:exp1a-prompt-results-staff} shows that the staff's evaluation confirmed the supervisor's key findings while being more conservative in their scoring approach across all criteria. Overall, the staff-based evaluation supports the fundamental trend observed in the supervisor's results: the common prompt consistently outperformed the complex prompt in nearly all quality-related criteria. The common prompt achieved superior staff ratings in \textit{Relevance} (9.50), \textit{Clarity} (9.18), \textit{Answerability} (9.43), \textit{Value} (7.26), \textit{Language} (9.27), and \textit{Correctness} (9.46), while the complex prompt only excelled in the \textit{Challenging} criterion (4.77 vs. 3.09).

\begin{table}[htbp]
    \centering
    \caption{Staff-based qualitative Analysis Results by Prompt Type (Exp1a).}
    \label{tab:exp1a-prompt-results-staff}
    \resizebox{\textwidth}{!}{%
    \begin{tabular}{lccccccccc}
        \toprule
        \textbf{Prompt} & \textbf{Metric} & \textbf{Relevance} & \textbf{Clarity} & \textbf{Answerability} & \textbf{Challenging} & \textbf{Value} & \textbf{Language} & \textbf{Correctness} & \textbf{Total Score} \\
        \midrule
        \multirow{3}{*}{Common} 
        & Mean & \textbf{9.50} & \textbf{9.18} & \textbf{9.43} & 3.09 & \textbf{7.26} & \textbf{9.27} & \textbf{9.46} & \textbf{57.19} \\
        & Std & 1.13 & 1.54 & 1.18 & 1.83 & 2.91 & 1.19 & 1.19 & 5.61 \\
        & Median & 10.0 & 10.0 & \textit{10.0} & 2.5 & \textit{8.0} & 10.0 & 10.0 & \textit{58.5} \\
        \midrule
        \multirow{3}{*}{Complex} 
        & Mean & 8.38 & 6.41 & 6.58 & \textbf{4.77} & 5.07 & 6.51 & 7.03 & 44.75 \\
        & Std & 2.35 & 2.95 & 3.30 & 2.53 & 3.19 & 2.88 & 3.17 & 12.11 \\
        & Median & 10.0 & 7.0 & 7.0 & \textit{5.0} & 5.0 & 7.0 & 8.0 & 44.5 \\
        \bottomrule
        \end{tabular}}
\end{table}

\pagebreak

However, the staff's stringent evaluation differs from the scores of the supervisor. For the common prompt, notable differences are visible in means of \textit{Clarity} (9.18 vs. 10.0), \textit{Value} (7.26 vs. 9.55), and \textit{Language} (9.27 vs. 10.0), while \textit{Correctness} remained remarkably similar between both groups (9.46 vs. 9.35). The lower ratings of the staff members are noticeable in the \textit{Total Score}, which shows that the common prompt was rated with a mean of 57.19 compared to the supervisor's 60.12.

Regarding the complex prompt, both evaluations consistently identified significantly lower performance across most criteria. The staff's ratings illustrate that this prompt type produces less relevant (8.38), clear (6.41), and answerable (6.58) questions. Most important for content fidelity, the complex prompt scored notably lower in \textit{Correctness} (7.03) compared to the common prompt (9.46), mirroring the supervisor's findings of 7.63 vs. 9.35.

The \textit{Total Score} difference between prompt types was noticeable across both evaluations. The staff showed a significant gap (57.19 vs. 44.75), while the supervisor demonstrated a similar one (60.12 vs. 47.25), with slightly higher ratings overall. This consistency strongly supports that the common prompt is significantly more effective for generating suitable questions, while the complex prompt -- producing more challenging questions -- compromises overall quality and correctness.

Overall, the staff-based evaluation confirmed the supervisor's key findings while being more stringent in their scoring. The complex prompt decreased the quality of the questions across most criteria, particularly in \textit{Clarity}, \textit{Answerability}, \textit{Value}, \textit{Language}, and \textit{Correctness}. In both evaluations, the \textit{Challenging} criterion was the only one where the complex prompt outperformed the common prompt. OpenAI also was consistently rated as the best-performing model across most criteria, while the input source preferences shifted a bit, with the staff members slightly rating the script-based questions as the best source material and the supervisor often favoring the transcript-based questions.

Besides these promising insights, the inter-rater agreement yielded undesirable results, as shown in Tables \bfref{tab:iaa-staff-1a} and \bfref{tab:iaa-combined-1a}. At first, Table \bfref{tab:iaa-staff-1a} shows the Fleiss' Kappa values which indicate poor agreement among the staff, with most criteria showing negative or low positive values. This suggests that the staff members had significant disagreements in their evaluations. This is due to the broad range of possible ratings (0-10) in the rubric, which could have lead to different interpretations of the criteria. 

\begin{table}[H]
   \centering
   \caption{Inter-Rater Agreement Analysis, only Staff (Exp1a).} 
   \label{tab:iaa-staff-1a}
   \begin{tabular}{lcccc}
        \toprule
        \textbf{Criterion} & \textbf{Fleiss' $\mathbf{\kappa}\uparrow$} & \textbf{Std} & \textbf{\#Items} & \textbf{\#Raters} \\
        \midrule
        Relevance      & -0.085 & 1.92 & 48 & 4 \\
        Clarity        &  0.034 & 2.72 & 48 & 4 \\
        Answerability  &  0.015 & 2.84 & 48 & 4 \\
        Challenging    &  0.011 & 2.35 & 48 & 4 \\
        Value          & -0.010 & 3.23 & 48 & 4 \\
        Language       &  0.001 & 2.59 & 48 & 4 \\
        Correctness    & -0.052 & 2.67 & 48 & 4 \\
        \bottomrule
   \end{tabular}
\end{table}

\pagebreak

The staff members' ratings were more diverse, while using all possible values in the rubric. Commonly, most scales were presented as 0, 2, 4, ..., 10, but the staff members also used values in between these steps, e.g.\ leading to one-off ratings.

In addition, the combined inter-rater agreement analysis in Table \bfref{tab:iaa-combined-1a} is presented, including the supervisor's ratings, solely using the desired values for the criteria, instead of the values in between. This table shows that the overall agreement across all raters was also poor, unfortunately. The combined agreement analysis led to slightly higher Fleiss' Kappa values for every criterion. 

\begin{table}[H]
   \centering
   \caption{Inter-Rater Agreement Analysis, combined (Exp1a).} 
   \label{tab:iaa-combined-1a}
   \begin{tabular}{lcccc}
        \toprule
        \textbf{Criterion} & \textbf{Fleiss' $\mathbf{\kappa}\uparrow$} & \textbf{Std} & \textbf{\#Items} & \textbf{\#Raters} \\
        \midrule
        Relevance      & -0.046 & 1.88 & 47 & 5 \\
        Clarity        &  0.077 & 2.68 & 47 & 5 \\
        Answerability  &  0.054 & 2.77 & 47 & 5 \\
        Challenging    &  0.014 & 2.61 & 47 & 5 \\
        Value          &  0.012 & 3.34 & 45 & 5 \\
        Language       &  0.043 & 2.53 & 48 & 5 \\
        Correctness    & -0.022 & 2.56 & 42 & 5 \\
        \bottomrule
   \end{tabular}
\end{table}

\object{Experiment 1b -- Error Handling}

As in the previous sub-experiment, the staff members were asked to evaluate the same experimental setup as the supervisor, shown in Table \bfref{tab:exp1b-manipulation-handling-staff}. This time, however, the focus was on error handling capabilities. Since the results are based on four staff members -- instead of one single supervisor --, the results are more diverse and provide a broader perspective on the \acp{llm}' performance in handling errors. The staff members evaluated the questions based on the same criteria as in Experiment 1a, but finally, the \textit{Manipulation Handling} criterion was solely regarded, since the other criteria were properly evaluated in the previous sub-experiment. 

\begin{table}[htbp]
    \centering
    \caption{Staff-based Manipulation Handling Analysis by LLM and Prompt Type (Exp1b).}
    \label{tab:exp1b-manipulation-handling-staff}
    \resizebox{.5\textwidth}{!}{%
    \begin{tabular}{llccc}
        \toprule
        & & \multicolumn{3}{c}{\textbf{Manipulation Handling}} \\
        \cmidrule{3-5}
        \textbf{LLM} & \textbf{Prompt Type} & \textbf{Mean} & \textbf{Std} & \textbf{Median} \\
        \midrule
        \multirow{2}{*}{Anthropic} 
        & Common & 2.50 & 4.63 & 0.0 \\
        & Complex & \textbf{7.25} & 2.55 & \textit{8.0}\\
        \midrule
        \multirow{2}{*}{DeepSeek} 
        & Common & 3.38 & 4.37 & 1.5 \\
        & Complex & \textbf{5.12} & 4.64 & \textit{5.5} \\
        \midrule
        \multirow{2}{*}{Google} 
        & Common & 3.50 & 4.66 & 0.5 \\
        & Complex & \textbf{7.88} & 3.00 & \textit{8.5} \\
        \midrule
        \multirow{2}{*}{OpenAI} 
        & Common & 2.50 & 4.63 & 0.0 \\
        & Complex & \textbf{5.00} & 3.93 & \textit{4.0} \\
        \bottomrule
    \end{tabular}}
\end{table}

\pagebreak

There are several key insights and differences compared to the supervisor's evaluation. Most notably, the staff members' scores were higher for \textit{Manipulation Handling}, especially visible in the complex prompt. Google and Anthropic achieved the best means of 7.88 and 7.25, while DeepSeek and OpenAI were rated 5.12 and 5.0. Regarding the common prompt, all models were rated lower, showing that this prompt type did not sufficiently encourage the models to detect manipulations.

Compared to the supervisor, who rated exclusively with the given values (0, 5, 10), the staff members used the full scale, rating all scores between 0 and 10. Here, Google and Anthropic were rated significantly higher for the complex prompt than by the supervisor (rated with 2.5), while OpenAI and DeepSeek still remained to be at the lower end, but with values $\geq 2.5$ instead. This stands in contrast to the quantitative analysis, where OpenAI and DeepSeek had shown better manipulation detection, based on rating every question instead of a sampled question set. Overall, the sample size remained small -- while the \textit{Adherence Score} of the quantitative analysis remains suitable --, but the qualitative evaluations still underscore the importance of prompt design. While the complex prompt improves \textit{Manipulation Handling} in Experiment 1b -- while being problematic in Experiment 1a via \textit{Correctness} --, the models still struggle to reliably identify and address manipulated content. 

The differences between the supervisor and staff ratings also illustrate the need for a more nuanced evaluation scale for \textit{Manipulation Handling}, as implemented in the other criteria. This would have been better with the \textit{Correctness} criterion in Experiment 1a as well. Unfortunately, the inter-rater agreement analysis Table \bfref{tab:iaa-staff-1b} also illustrates an unsatisfying agreement, compared to the previous sub-experiment. The Fleiss' $\kappa$ value between the staff members is 0.014, indicating just a slight agreement among these raters. This had happened because the staff members used the full range of scores instead of the desired ones, but overall, the ratings by all experts still present valuable information.

\begin{table}[H]
   \centering
   \caption{Inter-Rater Agreement Analysis (Exp1b).} 
   \label{tab:iaa-staff-1b}
   \begin{tabular}{lcccc}
        \toprule
        \textbf{Manipulation Handling} & \textbf{Fleiss' $\mathbf{\kappa}\uparrow$} & \textbf{Std} & \textbf{\#Items} & \textbf{\#Raters} \\
        \midrule
        Staff-only & 0.014 & 4.31 & 16 & 4 \\
        Combined with Supervisor & 0.036 & 4.21 & 15 & 5 \\
        \bottomrule
   \end{tabular}
\end{table}

Based on the staff members, there is no significant improvement compared to the previous sub-experiment, but the combined analysis -- including the supervisor -- shows a minimally higher Fleiss' $\kappa$ of 0.036, which is a slight value in the end. This is a little surprising, since the supervisor rated only the desired values (0, 5, 10), while the staff members used the full scale instead. The staff members also rated the questions with values $\geq 5$, which were not used by the supervisor, since the supervisor only rated up to 5 in certain cases. This individual misunderstanding in the criterion led to a low agreement between both groups, since the ratings were not comparable, even if the ideas behind the ratings were occasionally similar.

\pagebreak

\object{Mixed Analysis Complications} The qualitative evaluation revealed complications affecting both supervisor and staff assessments. Firstly, several questions were only partially or not at all related to the provided source material, limiting \textit{Answerability} and \textit{Correctness}. 

Few questions with unnecessary negations in them were criticized, reducing \textit{Clarity} and the \textit{Value}. This evaluation was further complicated by overly long or linguistically complex formulations, especially from Google and Anthropic using the complex prompt. Both groups often noted the use of generic phrases such as \enquote{based on the text}, which weakened the questions' educational \textit{Value}.

Factually incorrect or invented answer options -- which are not based on the source text -- were identified in multiple cases, making correctness difficult to assess, often seen at Google and Anthropic as well. Several questions were rated as very easy or trivial by both groups, particularly from OpenAI and DeepSeek, which is visible in the \textit{Challenging} criterion.

Questions based on the common prompt were often criticized for being too simple or not sufficiently challenging, while the complex prompt questions -- overly more challenging -- were often too convoluted or complex, leading to low values in e.g.\ \textit{Clarity} and \textit{Language}. Instead, the common prompt questions overall have shown higher ratings in both evaluations.

Furthermore, the strong preference of OpenAI in the supervisor's evaluation was not as pronounced in the staff members' ratings, but still slightly visible. The other models struggled more with achieving high scores in the criteria, especially in \textit{Clarity}, \textit{Answerability}, and \textit{Correctness}.

The evaluation rubric itself was not always clear for all criteria, resulting in inconsistent ratings and low inter-rater agreement among the groups. Overall, the staff members' ratings were more diverse, often using the full scale (0-10) instead of the discrete values (e.g.\ 0, 5, 10 at \textit{Correctness} and \textit{Manipulation Handling}) used by the supervisor. Moreover, the staff members rated the questions higher at \textit{Manipulation Handling} than the supervisor, indicating a different understanding of the criterion but still providing valuable insights based on more perspectives. 

In addition, the sampled question set in Experiment 1b was smaller than in Experiment 1a, which limits the representativeness of the results. At least, the staff-based qualitative analysis provided a broader perspective on the \acp{llm}' performance in handling errors, where the complex prompt improved \textit{Manipulation Handling} in Experiment 1b, while being problematic in Experiment 1a via \textit{Correctness}. This was not significantly visible in the supervisor's evaluation of Experiment 1b, but the quantitative analysis revealed similar trends according to the difference in prompt type performance with the \textit{Adherence Score}.

\subsection{Second Experiment}
\label{sec:second-experiment}

The second experiment addressed evaluating the relationship between question formats and Bloom's Taxonomy levels in \ac{llm}-generated questions. The goal was to determine how different question formats -- Multiple-Choice and Open-Ended questions -- perform across various cognitive levels of Bloom's revised Taxonomy, and also independently of each other. This experiment dealt with the second research question:

\requ{2}{How does the relationship between diverse question formats and Bloom's Taxonomy levels influence the pedagogical effectiveness of \ac{llm}-generated questions?}

\paragraph*{Qualitative Analysis -- Student-based} 

A group of three students, supervised by the secondary supervisor, evaluated the questions generated in Experiment 2. The students were asked to rate the questions based on the known criteria \textit{Relevance}, \textit{Clarity}, \textit{Answerability}, \textit{Challenging}, \textit{Value} and \textit{Language}. Additionally, the \textit{Bloom's Level} was assessed, as seen in Table \bfref{tab:exp2_criteria}. So, the evaluation process involved a detailed rubric, and the students' ratings were analyzed using Fleiss' Kappa again to assess inter-rater reliability.

Regarding Figure \bfref{fig:exp2-student-evaluation-metrics}, some key insights can be drawn from the student-based evaluation. First of all, the \textit{Relevance} criterion, showing that the questions were highly relevant to the topic across the models. The questions by Anthropic and DeepSeek were rated higher, achieving medians of 10 and means of 9.07 and 8.67, respectively. Google and OpenAI instead had lower ratings, with means of 8.2 and 8.0, coupled with a median of 8. Focusing on the \textit{Clarity} criterion, the models commonly achieved high scores, with all models showing a median of 10.0. The means were around 9, except for OpenAI, which had a mean of 8.67, but overall, the models generated logically clear questions.

Regarding \textit{Answerability}, all medians are 8. The means of the models slightly differ, showing that Anthropic and DeepSeek had higher means of 8.27 and 8.2, where DeepSeek shows a better data distribution. Focusing on Google and OpenAI, the outcomes are pretty similar, with the widest value ranges and means of 7.72 and 7.53, though still rated high. The \textit{Challenging} criterion presents medians of 8 for every model. Overall, more questions by Google and OpenAI were rated as more challenging, resulting in means of 7.27 and 7.47, respectively, while DeepSeek and Anthropic were rated lower with means of 6.8 and 6.53.

The results for the questions' \textit{Value} look pretty similar for each model, but the medians were slightly higher for DeepSeek and OpenAI (9.0) than for Google and Anthropic (8.0). Overall, the questions were rated as valuable, with minor outliers at DeepSeek and Google. Despite this, the means still remained high, with OpenAI achieving the highest mean of 8.73, followed by Anthropic (8.67), DeepSeek (8.53), and Google (8.33).

Furthermore, the \textit{Language} criterion looks similar across all models, with means from 8.37 (Google) to 8.73 (DeepSeek). In contrast, the \textit{Bloom's Level} shows a more diverse picture, where Google and OpenAI achieved the highest medians of 10 and 9.25, while DeepSeek and Anthropic were rated lower with medians of 6.0 and 7.0. So, many questions by Anthropic and DeepSeek were less rated at Bloom's revised Taxonomy or less aligned with the desired Bloom's level, if given.

Ultimately, the \textit{Total Score} illustrates that no \ac{llm} is a clear winner, with OpenAI being the lowest rated model with a mean of 55.27 and a median of 55.0 out of 70 points. The other models are rated slightly higher with lower ranges of values, but still showing that all models are quite similar in their performance, whether all criteria are considered together. For instance, a proper Multiple-Choice question by OpenAI from Experiment 2a demonstrates a positive example (Appendix \bfref{lst:2a-student-good-example}), while a problematic question by the same model shows complexity issues (Appendix \bfref{lst:2a-student-bad-example}).

\pagebreak

\begin{figure}[H]
   \centering
   \includegraphics[width=.9\textwidth]{../../../40_evaluation/exp2/qualitative/plots/exp2_boxplots_all_metrics_by_llm.png} 
   \caption{Student-based Evaluation Metrics by LLM (Exp2).}
   \label{fig:exp2-student-evaluation-metrics}
\end{figure}

\pagebreak

Regarding each sub-experiment instead, the results are shown in Figure \bfref{fig:exp2-student-evaluation-experiment}. Here, the three sub-experiments are compared, where Experiment 2a (Type Only) focused on solely specifying the question type, Experiment 2b focused on a specific Bloom level for each question (Bloom Only), and Experiment 2c combined both aspects (Type + Bloom). The \textit{Relevance} criterion shows that Experiment 2a achieved the highest mean of 9.04, followed by Experiment 2c with a mean of 8.12, and Experiment 2b with a mean of 8.08 and a wider range of values. Also, the median of Experiment 2a was the highest with 10, while Experiment 2b and 2c had medians of 8.

The \textit{Clarity} criterion shows a trend where Experiment 2c achieved the highest mean of 9.29, followed by Experiment 2a with a mean of 8.79, and Experiment 2b with 8.43, indicating that most of the questions were clear and understandable. Regarding the \textit{Answerability}, all sub-experiments achieved a median of 8. Experiment 2b was the worst rated one again, with a mean of 7.13, while Experiment 2a (mean of 8.17) and Experiment 2c (mean of 8.08) were rated better overall.

The \textit{Challenging} criterion shows an interesting trend, where Experiment 2b achieved the best results with a mean of 8.42, while Experiment 2a is following close behind with slightly shifted values and a mean of 7.17. The questions in Experiment 2c were rated as the least challenging with a mean of 6.17 instead. The questions' \textit{Value} was rated overly high across all sub-experiments, but the ones from Type Only (Experiment 2a) have a median of 8, while the other two sub-experiments achieved a median of 9. Overall, Experiment 2c had some outliers to the bottom, but overall, the means were still high, with 8.58 (Experiment 2a), 8.75 (Experiment 2b), and 8.46 (Experiment 2c).

In addition, the \textit{Language} criterion was rated overall best in Experiment 2c with a mean of 8.71, followed by Experiment 2a with a mean of 8.44, and Experiment 2b with a mean of 8.33. 2a and 2c achieved the highest medians of 10, while Experiment 2b had a median of 8.5 instead, which indicates that the questions were overall well-formed and linguistically precise, but the questions in 2b were slightly less linguistically valuable. Also, Experiment 2a share low scoring Experiment 2b, where the lowest scores are visible at 3, and Experiment 2c only goes down to 5 instead.

The \textit{Bloom Score} delivers a more interesting picture, where 2a is situated lower. In this sub-experiment, the mean is 4.0, showing that the questions are situated between \textit{Understanding} (3 points) and \textit{Applying} (4.5 points), since no Bloom level was specified in the prompt. This is acceptable, but the median is only 3.0 instead. In contrast, the values of 2b and 2c were situated significantly higher, with medians of 10 and means of 8.75 and 8.65, respectively. Overall, the Bloom levels were well adhered to (resulting in 10 points), with only a few outliers for both sub-experiments, illustrated via one-off ratings (5 points) or a rated Bloom level distance of $>1$ (0 points).

In this box plot, the \textit{Total Score} shows that the sub-experiments do not differ significantly, but Experiment 2b (Bloom Only) achieved the highest mean of 57.58 and a median of 58.5, followed by Experiment 2c (Type + Bloom) with a mean of 57.48 and a median of 58.5 as well. Experiment 2a (Type Only) is rated slightly lower with a mean of 54.19 and a median of 54.25, showing overall that the \textit{Bloom's Level} criterion is the most decisive one, since the scoring in Experiment 2a is significantly lower than in the other two sub-experiments. From this plot-based perspective, Experiment 2b and 2c are rated pretty similarly in \textit{Total Score}, while Experiment 2a also achieved acceptable results. 

\pagebreak

\begin{figure}[H]
   \centering
   \includegraphics[width=.9\textwidth]{../../../40_evaluation/exp2/qualitative/plots/exp2_boxplots_all_metrics_by_experiment.png}
   \caption{Student-based Evaluation Metrics by Sub-Experiment (Exp2).} 
   \label{fig:exp2-student-evaluation-experiment}
\end{figure}

\pagebreak

For instance, Anthropic illustrates a well-structured question from Experiment 2b for Bloom level 2 (Appendix \bfref{lst:2b-student-good-example}), while a problematic question from the same sub-experiment by Google demonstrates unnecessary complexity issues in its Bloom level 6 question, including many technical terms and a long introduction phrase (Appendix \bfref{lst:2b-student-bad-example}). Other examples from Experiment 2c by both models illustrate an inverted picture: one of Google's valuable \acp{mcq} for Bloom level 4 demonstrates a positive combination of question type and Bloom specification (Appendix \bfref{lst:2c-student-good-example}), while an Open-Ended question by Anthropic for the same Bloom level shows linguistic issues that can arise, based on complex wording and many foreign words (Appendix \bfref{lst:2c-student-bad-example}).

Since the \textit{Bloom's Level} criterion is the most important one in this experiment, Table \bfref{tab:exp2-llm-experiment-bloomslevel} shows the results of the Bloom levels achieved by the models in each sub-experiment. Here, certain trends among the models are visible, where Anthropic and Google achieved the highest means in Experiment 2a (4.96 and 4.54) -- which are situated at the level \textit{Applying} (4.5 points) --, while DeepSeek and OpenAI were rated even lower (2.83 and 3.67), which is close to \textit{Understanding} (3 points). In Experiment 2b, all models achieved the same mean of 8.33, except OpenAI, showing a perfect Bloom adherence with a mean of 10.0. Experiment 2c, combining both question type and Bloom level specification, shows that Google performed best with a mean of 9.58. Anthropic and DeepSeek follow with 8.75, while OpenAI is rated lowest with a mean of 7.50, still offering a good adherence to the Bloom level specification. The medians of Experiment 2b and 2c are rated 10/10 for every model, also indicating that the questions were well aligned with the specified Bloom levels. Experiment 2a instead shows that the medians were lower -- as this sub-experiment did not specify a Bloom level --, where Anthropic shows the highest median of 7, which is situated at \textit{Analyzing} (7 points). In contrast, the other models' questions were less cognitively demanding, illustrated by the median of 3.0 for each \ac{llm}, which is situated at the Bloom level \textit{Understanding} (3 points).

\begin{table}[htbp]
    \centering
    \caption[Student-based Bloom Scores by LLM and Sub-Experiment (Exp2).]{Student-based Bloom Scores by LLM and Sub-Experiment (Exp2).}
    \label{tab:exp2-llm-experiment-bloomslevel}
    \resizebox{.9\textwidth}{!}{%
    \begin{tabular}{lccccccccc}
        \toprule
        \multirow{2}{*}{\textbf{LLM}} & \multicolumn{3}{c}{\textbf{Mean}} & \multicolumn{3}{c}{\textbf{Std}} & \multicolumn{3}{c}{\textbf{Median}} \\
        \cmidrule(lr){2-4} \cmidrule(lr){5-7} \cmidrule(lr){8-10}
        & \textbf{Exp2a} & \textbf{Exp2b} & \textbf{Exp2c} & \textbf{Exp2a} & \textbf{Exp2b} & \textbf{Exp2c} & \textbf{Exp2a} & \textbf{Exp2b} & \textbf{Exp2c} \\
        \midrule
        Anthropic & \textbf{4.96} & 8.33 & 8.75 & 2.57 & 4.08 & 2.26 & \textit{7.0} & 10.0 & 10.0 \\
        DeepSeek  & 2.83 & 8.33 & 8.75 & 1.50 & 4.08 & 2.26 & 3.0 & 10.0 & 10.0 \\
        Google    & 4.54 & 8.33 & \textbf{9.58} & 2.53 & 4.08 & 1.44 & 3.0 & 10.0 & 10.0 \\
        OpenAI    & 3.67 & \textbf{10.00} & 7.50 & 2.78 & 0.00 & 3.99 & 3.0 & 10.0 & 10.0 \\
        \bottomrule
    \end{tabular}}
\end{table}

Since both question types (Multiple-Choice and Open-Ended) were equally generated in Experiments 2a and 2c, Table \bfref{tab:exp2-questiontype-experiment-bloomslevel} shows the results for each type. Overall, the Open-Ended questions were rated significantly higher than the \acp{mcq} in both sub-experiments. In Experiment 2a, Multiple-Choice achieved a mean of 2.42, rated between \textit{Remembering} (1.5 points) and \textit{Understanding} (3 points). Instead, Open-Ended questions achieved a mean of 5.58, which shows that questions of this type were more cognitively demanding -- between \textit{Applying} (4.5 points) and \textit{Analyzing} (7 points).

\pagebreak

Experiment 2c shows a similar trend, but focusing Bloom alignment instead of the cognitive level itself. The Bloom adherence for \acp{mcq} were rated high with a mean of 7.92. The Open-Ended questions were rated even higher with a mean of 9.38, showing a very high adherence to the specified Bloom level overall. The medians in Experiment 2c were 10.0 for both question types, but in Experiment 2a, the medians were 2.25 for Multiple-Choice and 7.0 for Open-Ended questions, indicating that the Open-Ended questions were more cognitively demanding, while \acp{mcq} were less cognitively demanding and less aligned with given Bloom levels. 

\begin{table}[htbp]
    \centering
    \caption[Student-based Bloom Scores by Question Format and each Sub-Experiment (Exp2).]{Student-based Bloom Scores by Question Format and each Sub-Experiment (Exp2).}
    \label{tab:exp2-questiontype-experiment-bloomslevel}
    \resizebox{.7\textwidth}{!}{%
    \begin{tabular}{lcccccc}
        \toprule
        \multirow{2}{*}{\textbf{Question Type}} & \multicolumn{2}{c}{\textbf{Mean}} & \multicolumn{2}{c}{\textbf{Std}} & \multicolumn{2}{c}{\textbf{Median}} \\
        \cmidrule(lr){2-3} \cmidrule(lr){4-5} \cmidrule(lr){6-7}
        & \textbf{Exp2a} & \textbf{Exp2c} & \textbf{Exp2a} & \textbf{Exp2c} & \textbf{Exp2a} & \textbf{Exp2c} \\
        \midrule
        Multiple-Choice & 2.42 & 7.92 & 1.23 & 3.27 & 2.25 & 10.0 \\
        Open-Ended      & \textbf{5.58} & \textbf{9.38} & 2.38 & 1.69 & \textit{7.00} & 10.0 \\
        \bottomrule
    \end{tabular}}
\end{table}

Regarding the inter-rater agreement analysis, the results are mixed, which is reflected in the varying Fleiss' kappa values across different criteria in Table \bfref{tab:iaa-students-exp2}. A major issue in the ratings was the diversity in scoring among the three students, indicating that there were differences in interpretation and application of the rubric criteria. Moreover, rating questions like the desired sample set is no everyday task for the students. This led to certain impacts on the agreement for \textit{Relevance}, \textit{Clarity}, \textit{Value}, and \textit{Language}. These criteria have a very low Fleiss' $\kappa$: \textit{Relevance} with 0.056, \textit{Clarity} with -0.061, \textit{Value} with -0.256, and \textit{Language} with 0.016.

\begin{table}[H]
    \centering
    \caption{Inter-Rater Agreement Analysis, Students (Exp2).}
    \label{tab:iaa-students-exp2}
    \begin{tabular}{lccc}
        \toprule
        \textbf{Criterion} & \textbf{Fleiss' $\mathbf{\kappa}\uparrow$} & \textbf{\#Items} & \textbf{\#Raters} \\
        \midrule
        Relevance          & 0.056 & 40 & 3 \\
        Clarity            & -0.058 & 40 & 3 \\
        Answerability      & 0.179 & 39 & 3 \\
        Challenging        & 0.270 & 40 & 3 \\
        Value              & -0.256 & 40 & 3 \\
        Language           & 0.016 & 40 & 3 \\
        Bloom Rating       & 0.397 & 40 & 3 \\
        \bottomrule
    \end{tabular}
\end{table}

In contrast to this, the \textit{Answerability} ($\kappa=0.179$) and especially the \textit{Challenging} ($\kappa=0.270$) criterion achieved better agreements. Moreover, the highest Kappa was achieved for the \textit{Bloom Rating} criterion ($\kappa=0.397$), indicating an almost moderate agreement. Here, the students' Bloom rating per question was compared, instead of taking the calculated \textit{Bloom's Level} score, based on the three scoring options in Table \bfref{tab:exp2_criteria}.

\pagebreak

In some cases, the students were not able to agree on the Bloom level of a question, which occurs since a set of questions did not provide a certain operator verb, or more than one operator verb was used, which both made it difficult to determine the final Bloom level of the question. Overall, the Kappa of 0.397 is still a good value for this criterion, showing that the students were able to agree on the Bloom levels in many cases, but still had some disagreements.

Overall, the stringent time constraints for the students' evaluation could have limited the depth of their analysis, as they had to evaluate 40 questions in a limited time frame. The unknown topic of the questions -- since these are pedagogical students -- may have also influenced the students' ability to accurately assess the questions' \textit{Relevance} and \textit{Answerability}, as they were not familiar with the ISO-OSI model before this analysis.

\object{Student-based Analysis Complications} This evaluation revealed further complications in the students' analysis affecting question quality across the sub-experiments. For instance, few questions without any operator verbs were present, particularly in \acp{mcq} of Experiment 2a, which are crucial for determining the Bloom level properly. This was visible in the \acp{mcq} across DeepSeek, Google, and OpenAI, where cognitive demands are hard to assess without clear operators. 

Questions by Anthropic, but particularly by OpenAI and Google across the sub-experiments were often criticized for excessive technical terminology and high linguistic complexity, making them overloaded and demotivating to read. Google often used complex wording, coupled with long introductory phrases, making these questions less valuable for educational purposes. Furthermore -- focusing Experiment 2a --, OpenAI was criticized for using too many operator verbs in the Open-Ended questions, which made it difficult to determine the \textit{Bloom's Level} and the \textit{Value} of the question, focusing educational usefulness.

Surprisingly, in Experiment 2c, one \textit{Creating}-level Multiple-Choice question (by OpenAI) was noted as unsuitable, since this question format constrained several responses, which contradicts this Bloom level. Based on difficulties like this, the students' ratings at \textit{Bloom's Level} were different from each other, which is reflected in the inter-rater agreement analysis in Table \bfref{tab:iaa-students-exp2}, but this criterion was still rated the highest by the students.

Overall, the mentioned issues were not present in every question, but these were visible in several questions across the models, which indicates that the \ac{llm}-generated questions were not always suitable for educational purposes. Also, the differences between the models were not as pronounced as in Experiment 1, where the models' performance was more diverse, which is an interesting finding.

Ultimately, one student criticized the \textit{Value} criterion, arguing it should be mapped to specific learning objectives rather than general educational worth, highlighting potential issues with the rubric design for students as evaluators who are unfamiliar with the technical domain.

\vp